% =====================================================================
%  Fiche d’activités : les 3 activités du chapitre 4 (physique)
%  Énergie potentielle de pesanteur et énergie mécanique — deux pages A4
%  (activités 1 et 2 page 1, activité 3 page 2)
%  (sans situation de départ, sans zone de nom, sans bandeau de titre)
%  À coller dans le cahier. Compilation : xelatex fiche_unique.tex (deux passes)
%  Le doc 4 est à l’échelle 1 : imprimer à 100 % (sans ajustement à la page).
% =====================================================================
\documentclass[10pt,a4paper]{article}
\usepackage[a4paper,top=.75cm,bottom=.6cm,left=.8cm,right=.8cm]{geometry}
\usepackage{fontspec}
\usepackage{polyglossia}
\setmainlanguage{french}
\usepackage{amsmath}
\usepackage{unicode-math}
\setmainfont{TeX Gyre Pagella}
\setsansfont{Poppins}[Path=fonts/, Extension=.ttf, UprightFont=*-Regular, BoldFont=*-Bold, ItalicFont=*-Italic, BoldItalicFont=*-BoldItalic]
\setmathfont{Latin Modern Math}
\setmathfont[range={it/{latin,Latin,greek,Greek}}]{texgyrepagella-italic.otf}
\setmathfont[range={up/{latin,Latin,num}}]{texgyrepagella-regular.otf}
\usepackage[table]{xcolor}
\usepackage{graphicx,tikz,array}
\usepackage[most]{tcolorbox}
\graphicspath{{images/}}
\usepackage{enumitem}
\usepackage{microtype}
\usepackage[newcommands]{ragged2e}
\setlength{\RaggedRightRightskip}{0pt plus 8em}
\usepackage[version=4]{mhchem}
\usepackage{siunitx}
\sisetup{locale=FR, per-mode=power, inter-unit-product=\ensuremath{\mathord{\cdot}}, group-minimum-digits=5}
\pagestyle{empty}
\setlength{\parindent}{0pt}
\setlength{\parskip}{0pt}
\relpenalty=10000 \binoppenalty=10000
\emergencystretch=3em

\definecolor{cInk}{HTML}{1B2A41}  \definecolor{cInk2}{HTML}{51627A}
\definecolor{cBlue}{HTML}{1F5FA8} \definecolor{cOrange}{HTML}{F07F1E}
\definecolor{cTeal}{HTML}{0E8F8A} \definecolor{cSoft}{HTML}{F4F7FB}
\definecolor{cRed}{HTML}{D62839}  \definecolor{cGreen}{HTML}{1F9D74}
\color{cInk}

\newcommand{\kw}[1]{{\sffamily\bfseries\color{cBlue}#1}}
\NewDocumentCommand{\docbadge}{m}{\tikz[baseline=(b.base)]\node[fill=cBlue,text=white,rounded corners=4pt,inner xsep=3pt,inner ysep=1.4pt,font=\sffamily\bfseries\tiny](b){Doc #1};}
% document : {largeur}{n}{légende}{image}
\newcommand{\doc}[4]{%
  \begin{minipage}[t]{#1}\vspace{0pt}\centering
    \includegraphics[width=\linewidth]{#4}\par\vspace{.5mm}
    {\sffamily\scriptsize\docbadge{#2}\enspace\color{cInk2}#3}
  \end{minipage}}
\NewDocumentCommand{\pastille}{m m}{\tikz[baseline=(n.base)]\node[circle,fill=#1,text=white,inner sep=0pt,minimum size=3.9mm,font=\sffamily\bfseries\tiny](n){#2};}
\setlist[enumerate,1]{label={\pastille{cBlue}{\arabic*}}, leftmargin=6mm, labelsep=1.8mm, itemsep=2pt, topsep=2pt, parsep=0pt}
% séparateur discret entre les activités
\newcommand{\sep}[1]{\par\vspace{2.5mm}%
  \noindent\tikz[baseline=(t.base)]\node[fill=cTeal,text=white,rounded corners=3pt,inner xsep=4pt,inner ysep=1.5pt,font=\sffamily\bfseries\scriptsize](t){Activité #1};%
  \hspace{2mm}{\color{cTeal!70}\leaders\hrule height .7pt\hfill\kern0pt}\par\vspace{1mm}}

% ---------- tableau de mesures compact ----------
\newcommand{\mesures}[3][0.8cm]{{\renewcommand{\arraystretch}{1.2}\setlength{\arrayrulewidth}{.5pt}\arrayrulecolor{cTeal!55}\setlength{\tabcolsep}{3pt}%
  \begin{tabular}{|>{\sffamily\scriptsize\bfseries\color{cTeal}}l|*{#2}{>{\centering\arraybackslash}p{#1}|}}\hline #3\end{tabular}}}

\newcommand{\vect}[1]{\overrightarrow{#1}}
\newcommand{\Ec}{E_{\mathrm{C}}}
\newcommand{\Epp}{E_{\mathrm{pp}}}
\newcommand{\Em}{E_{\mathrm{m}}}
\newcommand{\Eppi}[1]{E_{\mathrm{pp}#1}}
\newcommand{\Emi}[1]{E_{\mathrm{m}#1}}
\newcommand{\Eci}[1]{E_{\mathrm{C}#1}}
\newcommand{\WAB}[1]{W_{A\to B}(\vect{#1})}
\begin{document}
\small\RaggedRight
\setlength{\parskip}{0pt}
% =====================================================================
\sep{1}
\begin{minipage}[t]{0.69\linewidth}\vspace{0pt}\RaggedRight
\textbf{Doc 1.} Un solide (S) de masse $m$ se trouve au voisinage de la Terre, où le champ de pesanteur est uniforme, d’intensité $g$. On repère la position de son centre d’inertie $G$ par sa cote $z$ sur un axe $Oz$ vertical orienté vers le haut. Du fait de son interaction avec la Terre et de sa position, le solide possède une \textbf{énergie potentielle de pesanteur} (admise) : $\Epp = m\cdot g\cdot z + C$, où $C$ est une constante. Pour fixer $C$, on choisit une position $G_0$ d’altitude $z_0$, l’\textbf{état de référence}, où l’on pose $\Epp = 0$. Le solide passe de A (altitude $z_A$) à B (altitude $z_B$) ; la dénivellation est $h$.\par\vspace{.5mm}
{\color{cInk2}\textbf{Rappel.} $W_{A\to B}(\vect P) = m\cdot g\cdot(z_A - z_B)$ : moteur quand le solide descend, résistant quand il monte.}\par\vspace{.5mm}
\begin{enumerate}
  \item \kw{Définir.} Qu’est-ce que l’énergie potentielle de pesanteur d’un solide ? Quelle est son unité dans le Système international ?
  \item \kw{Déterminer.} Qu’est-ce que l’état de référence ? En écrivant que $\Epp = 0$ pour $z = z_0$, exprimer $C$ en fonction de $m$, $g$ et $z_0$.
  \item \kw{Établir.} En déduire $\Epp$ en fonction de $m$, $g$, $z$ et $z_0$, puis son expression lorsque $z_0 = 0$.
  \item \kw{Discuter.} Quel est le signe de $\Epp$ au-dessus et au-dessous de l’état de référence ? Comment varie $\Epp$ quand $z$ augmente ?
  \item \kw{Démontrer.} Exprimer $\Delta\Epp = \Eppi{B} - \Eppi{A}$ et démontrer la relation entre $\Delta\Epp$ et $\WAB P$.
\end{enumerate}
\end{minipage}\hfill
\doc{0.28\linewidth}{1}{Une caisse soulevée de A vers B}{doc1_repere}

% =====================================================================
\sep{2}
\begin{minipage}[t]{0.8\linewidth}\vspace{0pt}\RaggedRight
\textbf{Doc 2. Chute libre.} Un solide (S) en chute libre n’est soumis qu’à son poids $\vect P$. Il passe de A (altitude $z_A$, vitesse $V_A$) à B (altitude $z_B$, vitesse $V_B$).\par\vspace{.5mm}
\textbf{Doc 3. Plan incliné sans frottement.} Un mobile autoporteur de masse $m = \qty{732}{g}$ est abandonné sur une table à coussin d’air inclinée de $\alpha = 10^\circ$ ; frottements négligeables ; $g = \qty{9,8}{N.kg^{-1}}$. On enregistre les positions $M_1, \dots, M_8$ de son centre d’inertie toutes les $\tau = \qty{60}{ms}$ (doc 4). Référence de l’énergie potentielle : le plan horizontal passant par $M_8$. On note $d_i = M_{i-1}M_{i+1}$, $D_i = M_iM_8$ ; l’altitude de $M_i$ au-dessus de $M_8$ est $h_i = D_i\cdot\sin\alpha$.\par\vspace{.8mm}
\doc{0.5\linewidth}{3}{Table à coussin d’air inclinée}{doc3_table}\hfill
\begin{minipage}[t]{0.47\linewidth}\vspace{0pt}\RaggedRight
\begin{enumerate}
  \item \kw{Définir.} Écrire l’expression de l’énergie mécanique $\Em = \Ec + \Epp$ d’un solide en translation, de masse $m$ et de vitesse $V$, à l’altitude $z$.
  \item \kw{Démontrer.} Doc 2 : appliquer le théorème de l’énergie cinétique entre A et B, puis utiliser $\Delta\Epp = -\WAB P$ pour montrer que $\Emi{B} = \Emi{A}$. Que dire du poids ?
\end{enumerate}
\end{minipage}
\end{minipage}\hfill
\doc{0.17\linewidth}{2}{Chute libre de A à B}{doc2_chute}
\par\vspace{1mm}
{\centering\begin{minipage}{15.5cm}\centering
\includegraphics[width=15.5cm]{doc4_enregistrement}\par\vspace{.5mm}
{\sffamily\scriptsize\docbadge{4}\enspace\color{cInk2}Enregistrement à l’échelle 1 (positions du centre d’inertie) — \textbf{imprimer à 100 \%}}
\end{minipage}\par}
\vspace{1mm}
{\centering\mesures[0.95cm]{8}{%
Position & $M_1$ & $M_2$ & $M_3$ & $M_4$ & $M_5$ & $M_6$ & $M_7$ & $M_8$\\ \hline
$t$ (\unit{ms}) & 0 & 60 & 120 & 180 & 240 & 300 & 360 & 420\\ \hline
$d_i$ (\unit{cm}) & — & 1,30 & 2,30 & 3,50 & 4,60 & 5,90 & 7,10 & —\\ \hline
$D_i$ (\unit{cm}) & 14,5 & 14,0 & 13,2 & 11,7 & 9,7 & 7,1 & 3,8 & 0\\ \hline
$V_i$ (\unit{m.s^{-1}}) & & & & & & & & \\[1.2mm] \hline
$\Eci{i}$ (\unit{mJ}) & & & & & & & & \\[1.2mm] \hline
$\Eppi{i}$ (\unit{mJ}) & & & & & & & & \\[1.2mm] \hline
$\Emi{i}$ (\unit{mJ}) & & & & & & & & \\[1.2mm] \hline}\par}
\vspace{.5mm}
\begin{enumerate}[start=3]
  \item \kw{Calculer.} Docs 3 et 4 : exprimer $V_i$, $\Eci{i}$ et $\Eppi{i}$ du mobile en $M_i$ ; compléter le tableau, puis calculer $\Emi{i} = \Eci{i} + \Eppi{i}$. Que remarque-t-on ?
  \item \kw{Interpréter.} Tracer sur un même graphique $\Ec$, $\Epp$ et $\Em$ en fonction du temps. Interpréter l’allure des trois courbes.
  \item \kw{Formuler.} Le mobile est soumis à son poids et à la réaction normale $\vect R_N$ de la table. À quelle condition l’énergie mécanique d’un solide se conserve-t-elle ?
\end{enumerate}

\newpage
% =====================================================================
\sep{3}
\begin{minipage}[t]{0.69\linewidth}\vspace{0pt}\RaggedRight
\textbf{Doc 5. Débit d’air réduit.} Même dispositif que le doc 3, mais on réduit modérément le débit d’air : le mobile frotte un peu sur la table. On mesure sa vitesse $V$ et l’altitude $h$ de son centre d’inertie au-dessus de $P_5$ (état de référence). Entre $P_1$ et $P_5$, il parcourt $P_1P_5 = \qty{20,9}{cm}$.\par\vspace{.6mm}
{\centering\mesures[0.95cm]{5}{%
Position & $P_1$ & $P_2$ & $P_3$ & $P_4$ & $P_5$\\ \hline
$t$ (\unit{ms}) & 0 & 120 & 240 & 360 & 480\\ \hline
$V$ (\unit{m.s^{-1}}) & 0,100 & 0,268 & 0,436 & 0,604 & 0,773\\ \hline
$h$ (\unit{cm}) & 3,64 & 3,25 & 2,52 & 1,43 & 0\\ \hline
$\Ec$ (\unit{mJ}) & & & & & \\[1.2mm] \hline
$\Epp$ (\unit{mJ}) & & & & & \\[1.2mm] \hline
$\Em$ (\unit{mJ}) & & & & & \\[1.2mm] \hline}\par}\vspace{.6mm}
\textbf{Doc 6. Glissement avec frottement.} Un solide (S) de masse $m$ glisse de A (altitude $z_A$) vers B (altitude $z_B$) sur un plan incliné d’un angle $\alpha$. Il est soumis à son poids $\vect P$ et à la réaction $\vect R = \vect R_N + \vect f$ du plan : $\vect R_N$ est perpendiculaire au plan, $\vect f$ est une force de frottement parallèle au plan, opposée au mouvement, d’intensité constante $f$. $\vect R$, $\vect R_N$ et $\vect f$ sont appliquées au milieu I de la surface de contact.
\end{minipage}\hfill
\doc{0.28\linewidth}{6}{Glissement avec frottement de A vers B}{doc6_plan}
\par\vspace{.6mm}
\begin{minipage}[t]{0.49\linewidth}\vspace{0pt}\RaggedRight
\begin{enumerate}
  \item \kw{Exploiter.} Doc 5 : compléter le tableau. Comment évolue l’énergie mécanique du mobile ?
  \item \kw{Exprimer.} Doc 6 : exprimer $\Delta\Em = \Emi{B} - \Emi{A}$ en fonction de $\Delta\Ec$ et de $\Delta\Epp$.
  \item \kw{Établir.} Appliquer le théorème de l’énergie cinétique entre A et B, puis utiliser $\Delta\Epp = -\WAB P$ pour exprimer $\Delta\Em$ en fonction de $\WAB R$.
\end{enumerate}
\end{minipage}\hfill
\begin{minipage}[t]{0.49\linewidth}\vspace{0pt}\RaggedRight
\begin{enumerate}[start=4]
  \item \kw{Démontrer.} Avec $\vect R = \vect R_N + \vect f$, montrer que $\Delta\Em = \WAB f = -f\cdot AB$. Quel est son signe ?
  \item \kw{Interpréter.} Que devient l’énergie mécanique « perdue » ? Écrire la relation entre $\Delta\Em$, $\WAB f$ et la chaleur $Q$ dégagée.
  \item \kw{Estimer.} Doc 5 : en déduire la chaleur $Q$ dégagée entre $P_1$ et $P_5$ et l’intensité $f$ des frottements.
\end{enumerate}
\end{minipage}
\end{document}
